\documentclass{article}
\usepackage{amsmath,amssymb}
\usepackage{xurl}
\usepackage[hidelinks]{hyperref}
% Shared commands for notes in docs/paper-gaps/.

\DeclareUnicodeCharacter{2080}{\ifmmode{}_{0}\else\textsubscript{0}\fi}
\DeclareUnicodeCharacter{2081}{\ifmmode{}_{1}\else\textsubscript{1}\fi}
\DeclareUnicodeCharacter{2082}{\ifmmode{}_{2}\else\textsubscript{2}\fi}
\DeclareUnicodeCharacter{2099}{\ifmmode{}_{n}\else\textsubscript{n}\fi}

\newcommand{\Lean}{\textsc{Lean}}
\ExplSyntaxOn
\NewDocumentCommand{\leanid}{m}{
  \ifmmode
    \textsf{\detokenize{#1}}
  \else
    \group_begin:
    \urlstyle{sf}
    \tl_set:Nn \l_tmpa_tl {#1}
    \tl_replace_all:Nnn \l_tmpa_tl {\_} {_}
    \exp_args:NV \path \l_tmpa_tl
    \group_end:
  \fi
}
\ExplSyntaxOff
\providecommand{\tr}{\operatorname{tr}}
\newcommand{\tnleanIssueBase}{https://github.com/LionSR/TNLean/issues/}
\newcommand{\tnleanPrBase}{https://github.com/LionSR/TNLean/pull/}
\newcommand{\tnleanDocBase}{https://sirui-lu.com/TNLean/docs/}
\newcommand{\ghissue}[1]{\href{\tnleanIssueBase#1}{GitHub issue \##1}}
\newcommand{\ghpr}[1]{\href{\tnleanPrBase#1}{PR \##1}}
\newcommand{\doclink}[1]{\href{\tnleanDocBase#1.html}{\path{#1}}}

% Dirac notation and the identity operator, as in blueprint/src/macros/common.tex.
% \providecommand keeps note-local definitions authoritative.
\usepackage{dsfont}
\providecommand{\ket}[1]{|#1\rangle}
\providecommand{\bra}[1]{\langle #1|}
\providecommand{\braket}[2]{\langle #1 | #2 \rangle}
\providecommand{\ketbra}[2]{|#1\rangle\!\langle #2|}
\providecommand{\Id}{\mathds{1}}
\providecommand{\one}{\mathds{1}}
\providecommand{\C}{\mathbb{C}}
\providecommand{\MN}[1]{M_{#1}(\C)}
\providecommand{\MD}{M_D(\C)}

% Verdict marker: \gapnote{<kind>}{<status>}. Kind is the policy
% classification (clarification, local-correction, scope-restriction,
% unfaithful, false-source, open-gap); status is open, wip (elimination
% underway), resolved, or historical. Machine-read by the site index;
% prints a verdict line.
\newcommand{\gapnote}[2]{\par\noindent\textbf{Verdict:} #1 (#2).\par}

\title{Two-by-two regular PEPS renormalization with the original tensor}
\author{}
\date{2026-10-05}
\begin{document}
\maketitle
\gapnote{scope-restriction}{open}

\section{Source and exact result}
The source is Schuch, Cirac, and P\'erez-Garc\'ia,
\href{https://arxiv.org/abs/1001.3807}{arXiv:1001.3807},
Observations 6.5--6.6, local source lines 1818--1915 and the three diagrams
\path{figs4/renorm-gg-sym.pdf},
\path{figs4/renorm-g-id-and-split.pdf}, and
\path{figs4/renorm-fixedpoint.pdf}.

Let $G$ be any finite group, not necessarily abelian, and let $A$ be a
four-leg regular G-isometric tensor with any finite physical alphabet $P$.
On an untwisted doubled torus with coarse periods $w,h\geq3$, let $\Psi_f$
be the actual fine state and $\Psi_c$ the actual coarse state defined by
the same original tensor $A$. Both states are proved nonzero. There is an
isometry on the derived fine physical support, given by physical four-site
grouping, a product of local physical maps, and regrouping of output
registers, such that
\begin{align}
    I\Psi_f & =R\,\Psi_c\otimes\Omega,                 \\
    I\left(\frac{\Psi_f}{\|\Psi_f\|}\right)
            & =\frac{\Psi_c}{\|\Psi_c\|}\otimes\Omega,
    \qquad \|\Omega\|=1.
\end{align}
The exact positive scalar is
\begin{align}
    R=\left(\prod_{v\in\mathbb Z_w\times\mathbb Z_h}
    \frac{\sqrt{c_{B,v}}}{\sqrt{c_{C,v}}}\right)
    (\sqrt{|G|})^{|E(\Gamma)|}.
\end{align}
Here $c_{B,v}>0$ is the actual blocked tensor's isometry factor and
$c_{C,v}>0$ is the isometry factor of the original coarse tensor with its
surplus identity registers. All factors are derived from the local
G-isometry assumptions. The four-site Gram calculation separately gives
$c_{B,v}=|G|\prod_{i=0}^3c_{v,i}$ for the fine-site factors used in that
calculation; the group-order factor records the internal cycle.

\section{Geometric and physical content}
The fine torus is bijectively tiled by four-site blocks. An explicit
bijection identifies every fine horizontal and vertical bond with four
internal bonds and four paired crossing bonds per block. The eight separate
boundary labels are retained as the two labels on each coarse leg.
This derives the global fine-to-coarse contraction identity, without a
supplied factorization, Gram matrix, support condition, or nonzero-state
hypothesis.

The target is not merely a canonical tensor. Its local coefficient is the
original tensor on the reference labels times the equality indicator for
all relative surplus labels. Its G-isometry follows from the original
local Gram identity, and its graph contraction factors by an explicit
edge-label bijection. Thus the original-tensor comparison is a derived
physical operation. A separate theorem also gives an isometric equivalence
between the original tensor's full physical support and the canonical
averaging tensor's physical support.

All four original physical registers of each block are retained. The
physical maps act on the actual product of block ranges, pulled back to
the fine physical space by the full grouping unitary. The actual fine
state is proved to lie in this support. Unused ambient physical directions
are permitted, and neither ambient surjectivity nor an unspecified
extension to a unitary on those unused directions is assumed.

\section{Exact remaining scope}
The geometric bond-indexed reblocking is proved for every pair of positive
coarse periods, including one and two. The global Bell separation and
original-tensor normalized comparison use a simple coarse graph, so they
require both coarse periods at least three. Extending that physical
support argument to a bond-indexed graph with parallel bonds and loops is
still required for the smaller periods; no simple-graph theorem is applied
to them here.

The theorem recorded here has identity bonds. Its companion now proves
the corresponding uniform map for all native closures and coherent sums;
see \path{scp10_twisted_two_by_two_regular_fixed_point.tex} and
\leanid{TNLean.PEPS.exists_regularTwoByTwoTwistedTensorIsometry}. The common physical alphabet may be arbitrary,
and fine tensors may vary in the geometric and canonical-coarse results.
The normalized same-original-tensor statement assumes one homogeneous
four-leg tensor, as in the source's fixed-point assertion. No claim about
non-regular representations, approximate schemes, or unrestricted boundary
conditions is made. These limits keep the unrestricted source claim
separate from this proved periodic specialization.

\section{Declarations}
The exact geometric identity is
\leanid{TNLean.PEPS.torusBondNetwork_eq_twoByTwoBlocked}.
The derived four-site isometry is
\leanid{TNLean.PEPS.isGIsometric_twoByTwoSeparatedTensor}.
The canonical-coarse global support isometry is
\leanid{TNLean.PEPS.exists_regularTwoByTwoPhysicalSupportIsometry}.
The final original-tensor, scalar-correct and normalized statement is
\leanid{TNLean.PEPS.exists_regularTwoByTwoOriginalTensorIsometry}.
Nonvanishing is derived by the original identity closure sector, and the
factorization of the augmented original tensor is
\leanid{TNLean.PEPS.graphBondNetwork_regularGraphSurplusSite}.
\end{document}
